Tijdens die hete nacht
had ze seks niet verwacht.

Krekels staakten hun lied.
Menig scherpe grasspriet

kietelde op haar rug.
De maan trok zich terug

toen hij in haar kraamhok
met zijn boedel introk.

De donkere wolken
schoven als draaikolken

dreigend om elkaar heen,
terwijl soms een ster scheen

in de zwarte gaten.
De tuin leek verlaten

en de tuinman eenzaam.
Ze gaf hem haar lichaam.

Hij gaf haar zijn vruchten.
Wat had ze te duchten?

Ze had naar hem gesmacht,
maar bleek ineens verkracht

nadat ze aanpapte
en men hen betrapte.

Dat werd haar verklaring.
Er volgde ophanging

als straf voor haar vrijer.
Ze werd zo van vleier

en slinkse uitdager
een valse aanklager

en gladde meeprater.
Toch sloeg ze een flater

toen de straf werd gevierd
en het lijk werd ontsierd,

bij de officiaal.
Ook veinsde ze ditmaal

inlevingsvermogen
bij haar dialogen

puur uit eigen gewin
met haar standaard vangzin.

“Heb jij daar ook last van?”,
werkte niet bij een man.

De knecht van de bisschop
had het over zijn kop.

Hij verloor zijn haren,
schaarser door de jaren

dat hij recht moest spreken.
Haar kans leek verkeken.

Haar schijn, te laat doorzien,
kon niet worden herzien.

Religie was voor haar
niet per se altijd waar.

Die set van convicties
volgde ze wel precies,

maar uit conformisme
of opportunisme

om het goede te doen.
Haar brevet van fatsoen.

Een leven als een sloof.
Ze hoopte haar geloof,

aldus te verliezen,
zodat ze kon kiezen

voor een open leven,
dat niet werd beschreven

in rare geschriften
van oude geschiften.

Ze wilde ontdekken
hoe die vreemde gekken

ooit overtuigingen,
die standaard rondgingen

plots naast zich neerlegden
en iets waarzegden.

Net als deze saga
met de rol "Maria".

Het boek van het leven
zelf opnieuw omschreven.

De bijbel was zo’n boek.
Voor haar gesneden koek

was de aard van het werk.
Het versterkte de kerk,

maar was niet historisch
hoogstens filosofisch.

Een overlevering,
die vaak slechts mondeling

ooit was opgetekend
en dus een vertekend

beeld aan ons kon geven.
Het was vaak een streven

van een alleenheerser
om elke hardleerser

toch te overtuigen
zonder af te tuigen.

Hoewel laatstgenoemde
uiteraard opdoemde

als religie faalde
en een rebel smaalde.
In al die verhalen
werd bij het vertalen

ook vaak fouten gemaakt.
Haar cel werd niet gekraakt

door simpel wat woorden,
die bovendien smoorden

jegens de ervaring,
de echte beleving.

Dit leidde tot argwaan.
Met name trok haar aan

wat de kerk haar verbood
De weidse naamgenoot,

van de non Maria,
genaamd Magdalena,

werd door de kerk geweerd.
En Maria had geleerd

dat deze volgeling
van Jezus als leerling

niets goeds te melden had.
Maria bestreed dat.

In de rust van haar cel
mocht Magdalena wel

hangen aan haar celdeur.
Daar doorbrak ze haar sleur.

In plaats van heiligheid
kreeg genot zijn vrijheid.

Mocht ze zich wel geven
aan haar gevoelsleven.

Kreeg lust als thema
in Emma’s opera
een gelaagd karakter.
Werd de zucht exacter

van de deugd gescheiden.
Kon haar zonde leiden

van schuld tot genade
vanuit de façade,

die paus Gregorius
op de “hoer” van Jezus

uit vrouwenhaat plakte.
Want Maria snakte

om te zijn als de vrouw,
die Jezus als lief wou.

“Lieve Magdalena
uw dienaar Maria

vraagt om uw genade
voor mijn escapade.

Ook ik gaf me over
en loog zelfs daarover

aan de kracht van genot.
Echter de zoon van God

dreef zonder een geluid,
die zonde bij u uit.

En nog zes anderen.
U mocht veranderen

van een zondig wezen
naar door God verrezen

zonder hard te werken.
Kunt u mij niet sterken

hier biddend in mijn cel.
U, als de apostel

der apostelen, kunt
wat mij niet is gegund,

vast en zeker geven.
Geef mij een heilig leven!”

Er was een boek edoch
namelijk van Enoch,

die Maria aannam
alsof het van God kwam

omdat haar eigen kerk
dit oud religieus werk

uit de bijbel weerde.
Dat terwijl ze leerde

van juist de kerkvaders
de bredere kaders

van dit apocrief boek.
Ze ging daarom op zoek

naar alle citaten
die wellicht de gaten

voor haar konden dichten.
De kern ging om plichten

die niet werden vervuld.
Het engelengeduld

werd op de proef gesteld.
In het aardse werkveld

kwam men mensen tegen.
De engelen kregen

lustgevoelens voor hen.
Dit leidde tot seks en

tot een ras van reuzen.
God strafte de kneuzen

in Henochs visioen.
Emmy besefte toen

al meelezend via
de blik van Maria,

wat moeder daarop zei:
“De reuzen dat zijn wij!

Wij zijn de Nephilim.
Zo sprak onze pelgrim,

de heilige profeet.”
Emmy was niet gereed

voor deze ontdekking.
Het bewees de strekking

van haar dwaze moeder.
De reus had als broeder

een engel of wachter.
Eén ouder was zachter

dan de ruigere mens.
Lichtvoetig maar intens.

Vlug van tong, ver van huis,
rein maar blijkbaar onkuis.

De lijfbehoudzwoeger
vluchtte eerst naar vroeger.

Naar haar zuster Emma.
Dat bleek een dilemma

want zij omarmde lust
volkomen en bewust.

Wie had anders verwacht?
De wet van het geslacht!

Alles is vrouwelijk
als zowel mannelijk

wat betreft energie.
Dat ondervond Emmy.

Het duurde maar even.
Want in een vorig leven

zocht Emmy toen haar heil.
Zelfs Maria werd geil

en gaf haar geen redding.
Maar Jezus’ volgeling

genaamd Magdalena
kwam in haar arena

als fijne paria
via haar Maria.

Wensbeelden vielen om.
Leeg en zoekend alom

berokkende haar schrik.
Met alleen haar rest-ik

kolkte Emmy terug.
Beheerst doch vliegensvlug.
